\chapter{Sprays and geodesics}
\label{chap:geodesics}

Sprays, the correspondence between quadratic sprays and symmetric connections, and geodesics on
Banach manifolds: the geodesic equation in charts, local existence with uniform time, uniqueness,
maximal geodesics, the exponential map and normal neighbourhoods, after Lang
\cite[\S\S IV.3--4, VIII.5]{Lan99}. In general locally convex spaces there is no existence theorem
for ordinary differential equations, so existence is only claimed for Banach models.

\section{Sprays}

A spray, or second order vector field, is a vector field $\xi$ on $TM$ such that
$T\pi \circ \xi = \id$. In a chart it has the form $(x, v) \mapsto (x, v; v, f(x, v))$, and it is
determined by its \emph{local part} $f$. As we did for connections, we encode sprays through their
local parts.

\begin{definition}[Spray]
  \label{def:spray}
  \lean{LocallyConvex.Manifolds.Spray, LocallyConvex.Manifolds.Spray.IsQuadratic}
  \leanok
  \uses{def:mb-manifold, def:tangent-bundle}
  A \emph{spray} on $M$ assigns to each chart $i$ of the atlas a smooth map
  $f_i \colon i(U_i) \times E \to E$ such that, for every pair of charts $i, j$, every $y$ in the
  domain of $\tau = j \circ i^{-1}$ and every $v \in E$,
  \[
    f_j\bigl(\tau y, \dd\tau(y)(v)\bigr) = \dd\tau(y)\bigl(f_i(y, v)\bigr) + \dd^{(2)}\tau(y; v, v).
  \]
  The spray is \emph{quadratic} if $f_i(x, s v) = s^2 f_i(x, v)$ for all $i$, $x$, $v$ and
  $s \in \R$.
\end{definition}

The change of charts law is the change of charts law of the double tangent bundle $T(TM)$,
restricted to the points of the form $(x, v; v, w)$. In \cite{Lan99}, \S IV.3, a spray is by
definition quadratic; we keep the two notions apart, since the quadratic condition is what is
needed to pass to connections.

\begin{proposition}[The spray of a connection]
  \label{prop:conn-to-spray}
  \lean{LocallyConvex.Manifolds.Connection.toSpray, LocallyConvex.Manifolds.Connection.toSpray_isQuadratic}
  \leanok
  \uses{def:connection, def:spray}
  Every connection defines a quadratic spray, with local parts $f_i(x, v) = \Gamma_i(x)(v, v)$.
\end{proposition}

\begin{proof}
  \leanok
  Restrict the change of charts law of $\Gamma$ to the diagonal $a = b = v$. Homogeneity of degree
  two follows from bilinearity.
\end{proof}

The converse requires recovering a bilinear map from its quadratic form, and checking that the
polarization obeys the change of charts law of connections. For the first point, the following
identity holds without any regularity assumption.

\begin{lemma}[Homogeneous maps of degree two]
  \label{lem:conn-homogeneous}
  \lean{LocallyConvex.Manifolds.hasDirDerivAt_of_homogeneous, LocallyConvex.Manifolds.dMB_dMB_zero_of_homogeneous, LocallyConvex.Manifolds.polarization}
  \leanok
  Let $g \colon E \to F$ satisfy $g(s v) = s^2 g(v)$ for all $s \in \R$, $v \in E$. Then
  $\dd^{(2)} g(0; v, v) = 2\, g(v)$. Moreover, if $B$ is symmetric and bilinear, then
  $B(a + b, a + b) - B(a - b, a - b) = 4 B(a, b)$.
\end{lemma}

\begin{proof}
  \leanok
  By homogeneity, $g(s v + t v) - g(s v) = \bigl((s + t)^2 - s^2\bigr) g(v)$, so the difference
  quotient can be computed exactly and $\dd g(s v)(v) = 2 s\, g(v)$ for every $s$.
  Differentiating once more along the line $s \mapsto s v$ at $s = 0$ gives $2 g(v)$. The
  polarization identity is immediate.
\end{proof}

\begin{theorem}[The connection of a quadratic spray]
  \label{thm:conn-spray-to-connection}
  \lean{LocallyConvex.Manifolds.Spray.polar, LocallyConvex.Manifolds.Spray.polar_eq, LocallyConvex.Manifolds.Spray.polar_transition, LocallyConvex.Manifolds.Spray.toConnection, LocallyConvex.Manifolds.Spray.toConnection_isSymmetric, LocallyConvex.Manifolds.Spray.toConnection_toSpray}
  \leanok
  \uses{def:spray, def:connection, lem:conn-homogeneous, prop:conn-to-spray}
  Let $f$ be a quadratic spray. In each chart $i$ put
  \[
    \Gamma_i(x)(a, b) = \tfrac14\bigl(f_i(x, a + b) - f_i(x, a - b)\bigr)
      = \tfrac12\, \dd^{(2)}\bigl(f_i(x, \cdot)\bigr)(0; a, b).
  \]
  Then the $\Gamma_i$ are the symbols of a symmetric connection, whose spray is $f$.
\end{theorem}

\begin{proof}
  \leanok
  By Lemma~\ref{lem:conn-homogeneous} applied to $g = f_i(x, \cdot)$ and to the symmetric
  bilinear map $\tfrac12 \dd^{(2)} g(0; \cdot, \cdot)$, the two expressions agree; the second one
  is bilinear and symmetric, and $\Gamma_i(x)(v, v) = f_i(x, v)$. Smoothness of the symbols
  follows from the first expression. The change of charts law for $\Gamma$ is obtained by
  polarizing the change of charts law of the spray, using that $\dd\tau(y)$ is linear and
  $\dd^{(2)}\tau(y; \cdot, \cdot)$ is symmetric and bilinear. This is \cite{Lan99}, Prop.~IV.3.2
  and \S VIII.2.
\end{proof}

Thus quadratic sprays and symmetric connections are the same object, through
$f(x, v) = \Gamma_x(v, v)$; in the rest of the chapter we work with connections.

\section{Geodesics on Banach manifolds}

In this section $E$ is a Banach space. Then Michal--Bastiani and Fr\'echet derivatives agree, and
in a chart $i$ the geodesic equation is the second order equation
\[
  c'' = \Gamma_i(c)(c', c')
\]
on the open set $\varphi_i(U_i) \subseteq E$, with $\Gamma_i$ smooth and bilinear in the last two
variables. We say that $\alpha$ is a \emph{geodesic on an open set} $I \subseteq \mathbb{R}$ if it is
continuous on $I$ and satisfies this equation at every $t \in I$ in the preferred chart at
$\alpha(t)$.

\begin{theorem}[The geodesic equation does not depend on the chart]
  \label{thm:geodesic-eq-chart}
  \lean{LocallyConvex.Manifolds.geodesicEq_comp, LocallyConvex.Manifolds.Connection.geodesicEqAt_of_chart}
  \leanok
  If $\alpha$ is continuous at $t$ and satisfies the geodesic equation at $t$ in a chart around
  $\alpha(t)$, it satisfies it in every chart around $\alpha(t)$.
\end{theorem}

\begin{proof}
  \leanok
  With $\tau = \varphi_j \circ \varphi_i^{-1}$ and $c = \varphi_i \circ \alpha$,
  $(\tau \circ c)'' = D^2\tau(c', c') + D\tau(c'')$, and the change of charts law of the symbols
  reads $\Gamma_j(\tau y)(D\tau\, a, D\tau\, b) = D\tau\, \Gamma_i(y)(a, b) + D^2\tau(a, b)$.
\end{proof}

\begin{theorem}[Local existence of geodesics, with uniform time]
  \label{thm:geodesic-existence}
  \lean{LocallyConvex.Manifolds.Connection.exists_isGeodesicOn_uniform, LocallyConvex.Manifolds.Connection.exists_isGeodesicOn}
  \leanok
  \uses{thm:geodesic-eq-chart}
  Let $y_0$ be a point of the chart $i$. There is $r > 0$ such that for every $K \ge 0$ there is
  $\varepsilon > 0$ with the following property: for every $y$ with $\|y - y_0\| \le r$ and every
  $w$ with $\|w\| \le K$ there is a geodesic $\alpha$ on $(-\varepsilon, \varepsilon)$, staying in
  the domain of $i$, with $\varphi_i(\alpha(0)) = y$ and $(\varphi_i \circ \alpha)'(0) = w$.
\end{theorem}

\begin{proof}
  \leanok
  Picard--Lindel\"of for the first order system $(c, v)' = (v, \Gamma_i(c)(v, v))$, whose right
  hand side is $\Ck{1}$, hence locally Lipschitz; the time is uniform for small velocities, and the
  homogeneity of the equation ($t \mapsto c(st)$ is a solution with initial velocity $s\, c'(0)$)
  makes it uniform for velocities bounded by $K$.
\end{proof}

\begin{theorem}[Uniqueness of geodesics]
  \label{thm:geodesic-uniqueness}
  \lean{LocallyConvex.Manifolds.Connection.IsGeodesicOn.eqOn}
  \leanok
  \uses{thm:geodesic-eq-chart}
  If $M$ is Hausdorff, two geodesics on $(a, b)$ with the same position and velocity at some
  $t_0 \in (a, b)$ agree on $(a, b)$.
\end{theorem}

\begin{proof}
  \leanok
  The set of $t$ near which the two geodesics agree is open. It is also closed in $(a, b)$: at a
  limit point $t$ the two curves meet, because $M$ is Hausdorff; in a chart around this common
  point both solve the same equation, which has unique solutions since its right hand side is
  locally Lipschitz. By connectedness it is all of $(a, b)$.
\end{proof}

\begin{remark}
  The Hausdorff hypothesis cannot be dropped: on the line with two origins, the two copies of a
  line through the origins are geodesics of the flat connection with the same data at $t = 1$, and
  they differ at $t = 0$.
\end{remark}

\begin{theorem}[Extension of geodesics]
  \label{thm:geodesic-extension}
  \lean{LocallyConvex.Manifolds.Connection.IsGeodesicOn.exists_extend}
  \leanok
  \uses{thm:geodesic-existence, thm:geodesic-uniqueness}
  Let $M$ be Hausdorff, and let $\alpha$ be a geodesic on $(a, b)$ such that $\alpha(t) \to q$ as
  $t \to b^-$ and the velocity of $\alpha$, read in a chart around $q$, stays bounded as
  $t \to b^-$. Then $\alpha$ extends to a geodesic on $(a, b')$ for some $b' > b$.
\end{theorem}

\begin{proof}
  \leanok
  Let $\varepsilon$ be the uniform time of Theorem~\ref{thm:geodesic-existence} around
  $\varphi_i(q)$. Take $t_1 > b - \varepsilon/2$ with $\alpha(t_1)$ close to $q$, and the geodesic
  $\mu$ with the position and velocity of $\alpha$ at $t_1$. By uniqueness
  $\mu(t - t_1) = \alpha(t)$ where both are defined, so the two curves glue to a geodesic on
  $(a, t_1 + \varepsilon)$, and $t_1 + \varepsilon > b$.
\end{proof}

\begin{theorem}[Complete Finsler manifolds: geodesics extend {\cite[Prop.~9.2.1]{Lar22}}]
  \label{thm:finsler-geodesic-extension}
  \lean{LocallyConvex.Manifolds.Connection.IsGeodesicOn.exists_extend_of_complete}
  \leanok
  \uses{thm:geodesic-extension}
  Let $M$ be Hausdorff, with a Finsler length structure $\mu$ bounded below, such that $M$ is
  complete for the symmetrized distance $\max(d_b(x, y), d_b(y, x))$. If $\alpha$ is a geodesic on
  $(a, b)$ with $\mu(\dot\alpha)$ and $\mu(-\dot\alpha)$ bounded, then $\alpha$ extends to a
  geodesic on an interval $(a', b') \supsetneq (a, b)$ at both ends.
\end{theorem}

\begin{proof}
  \leanok
  Reparametrizing arcs of $\alpha$ by a smooth ramp gives $\Ck{1}$ competitors for $d_b$, so
  $d_b(\alpha(s), \alpha(t)) \le L(t - s)$ and the same backwards; thus $\alpha(t)$ is Cauchy as
  $t \to b^-$. Its limit in $(M, d_b)$ is a limit in $M$ because $\mu$ is bounded below, and the
  lower bound also bounds the velocity in a chart around the limit point.
  Theorem~\ref{thm:geodesic-extension} applies; the end $a$ is treated by reversing time.
\end{proof}

\begin{remark}
  In \cite{Lar22} the metric is invariant under the parallel transport of the spray, so geodesics
  have constant speed; we only use bounded speed. Since $\mu$ need not be reversible, both
  $\mu(\dot\alpha)$ and $\mu(-\dot\alpha)$ have to be bounded. With maximal geodesics, the
  statement ``every geodesic is defined on $\mathbb{R}$'' is Theorem~\ref{thm:finsler-geodesically-complete}.
\end{remark}

\begin{theorem}[Maximal geodesics {\cite[\S IV.3, \S IV.4]{Lan99}}]
  \label{thm:maximal-geodesic}
  \lean{LocallyConvex.Manifolds.tangentSMul,
    LocallyConvex.Manifolds.Connection.IsGeodesicOn.eqOn_of_ordConnected,
    LocallyConvex.Manifolds.Connection.IsGeodesicThrough,
    LocallyConvex.Manifolds.Connection.exists_isGeodesicThrough,
    LocallyConvex.Manifolds.Connection.maxInterval,
    LocallyConvex.Manifolds.Connection.maxGeodesic,
    LocallyConvex.Manifolds.Connection.isOpen_maxInterval,
    LocallyConvex.Manifolds.Connection.isGeodesicThrough_maxGeodesic,
    LocallyConvex.Manifolds.Connection.IsGeodesicThrough.subset_maxInterval,
    LocallyConvex.Manifolds.Connection.IsGeodesicThrough.maxGeodesic_eq,
    LocallyConvex.Manifolds.Connection.maxInterval_smul,
    LocallyConvex.Manifolds.Connection.maxGeodesic_smul,
    LocallyConvex.Manifolds.Connection.maxInterval_zero_smul}
  \leanok
  \uses{thm:geodesic-existence, thm:geodesic-uniqueness}
  Let $M$ be Hausdorff and modeled on a Banach space, with a connection. The geodesics with
  initial velocity $v \in \T M$ glue to a \emph{maximal geodesic} $\gamma_v$, defined on an open
  interval $J_v \ni 0$: every geodesic with initial velocity $v$ on an open interval $I \ni 0$ is
  the restriction of $\gamma_v$ to $I$. For $s \neq 0$, $J_{sv} = s^{-1}J_v$ and
  $\gamma_{sv}(t) = \gamma_v(st)$; the maximal geodesic with zero initial velocity is constant and
  defined on $\mathbb{R}$.
\end{theorem}

\begin{proof}
  \leanok
  As for the maximal flow (Theorem~\ref{thm:maximal-flow}), from local existence and uniqueness on
  intervals; homogeneity because $t \mapsto \alpha(st)$ is a geodesic with initial velocity
  $s\,\dot\alpha(0)$. The maximal geodesic is built directly from the geodesics, not as the
  projection of the maximal flow of the geodesic spray on $\T M$; the two agree by
  Theorem~\ref{thm:geodesic-spray-flow}. We do not claim this argument is new.
\end{proof}

\begin{definition}[The exponential map {\cite[\S IV.4]{Lan99}, \cite[\S9.2]{Lar22}}]
  \label{def:exp}
  \lean{LocallyConvex.Manifolds.Connection.expDomain, LocallyConvex.Manifolds.Connection.exp,
    LocallyConvex.Manifolds.Connection.expDomainAt, LocallyConvex.Manifolds.Connection.expAt,
    LocallyConvex.Manifolds.Connection.smul_mem_expDomain_iff,
    LocallyConvex.Manifolds.Connection.exp_smul}
  \leanok
  \uses{thm:maximal-geodesic}
  The \emph{exponential map} is $\exp v = \gamma_v(1)$, on $D(\exp) = \{v : 1 \in J_v\}$, and
  $\exp_p$ is its restriction to $T_pM$. Then $tv \in D(\exp)$ if and only if $t \in J_v$, and
  $\exp(tv) = \gamma_v(t)$; so the maximal geodesic is $t \mapsto \exp(tv)$.
\end{definition}

\begin{definition}[The geodesic spray {\cite[\S IV.3, pp.~97--98]{Lan99}}]
  \label{def:geodesic-spray}
  \lean{LocallyConvex.Manifolds.Connection.sprayFun,
    LocallyConvex.Manifolds.Connection.ddChart_sprayFun,
    LocallyConvex.Manifolds.Connection.spray,
    LocallyConvex.Manifolds.Connection.rep_spray,
    LocallyConvex.Manifolds.Connection.spray_isSecondOrder}
  \leanok
  \uses{def:connection}
  Let $M$ be modeled on a Banach space, with a connection. The \emph{geodesic spray} is the
  smooth vector field $S$ on the manifold $\T M$ (modeled on $E \times E$) which, in the chart of
  $\T M$ over a chart $i$ of $M$, reads
  \[
    (x, w) \longmapsto (x, w;\ w, \Gamma_i(x)(w, w)).
  \]
  It is well defined: under the transition maps of $\T\T M$, the local formulas in two charts
  agree by the change of charts law of the symbols. It is a second-order vector field,
  $\pi_* S(v) = v$.
\end{definition}

\begin{theorem}[Geodesics and the flow of the spray {\cite[\S IV.3, p.~97]{Lan99}}]
  \label{thm:geodesic-spray-flow}
  \lean{LocallyConvex.Manifolds.Connection.IsGeodesicOn.isIntegralCurveOn_velocity,
    LocallyConvex.Manifolds.Connection.isGeodesicOn_proj_of_isIntegralCurveOn,
    LocallyConvex.Manifolds.Connection.velocity_proj_of_isIntegralCurveOn,
    LocallyConvex.Manifolds.Connection.maxInterval_eq_spray,
    LocallyConvex.Manifolds.Connection.maxGeodesic_eq_proj_flow,
    LocallyConvex.Manifolds.Connection.flow_spray_eq_velocity}
  \leanok
  \uses{def:geodesic-spray, thm:maximal-geodesic, thm:maximal-flow}
  The velocity $\dot\alpha$ of a geodesic is an integral curve of $S$; conversely, an integral
  curve $\beta$ of $S$ is the velocity of its projection $\pi\beta$, which is a geodesic. Hence,
  if $M$ is Hausdorff, the maximal interval $J_v$ of the geodesic with initial velocity $v$ is the
  maximal interval of the integral curve of $S$ through $v$, $\gamma_v(t) = \pi(\mathrm{Fl}_S(t, v))$
  and $\mathrm{Fl}_S(t, v) = \dot\gamma_v(t)$.
\end{theorem}

\begin{proof}
  \leanok
  In the chart of $\T M$ over a chart $i$ of $M$, an integral curve of $S$ is a curve $(x, w)$
  with $x' = w$ and $w' = \Gamma_i(x)(w, w)$, which is the geodesic equation for $x$.
\end{proof}

\begin{theorem}[$D(\exp)$ is open and $\exp$ is smooth {\cite[\S IV.4, p.~107]{Lan99},
  \cite[\S6.1.2]{Lar22}}]
  \label{thm:exp-smooth}
  \lean{LocallyConvex.Manifolds.Connection.expDomain_eq_preimage_flowDomain,
    LocallyConvex.Manifolds.Connection.exp_eq_proj_flow,
    LocallyConvex.Manifolds.Connection.isOpen_expDomain,
    LocallyConvex.Manifolds.Connection.smoothMBMOn_exp,
    LocallyConvex.Manifolds.Connection.isOpen_expDomainAt,
    LocallyConvex.Manifolds.Connection.smoothMBMOn_expAt,
    LocallyConvex.Manifolds.Connection.zero_mem_expDomain,
    LocallyConvex.Manifolds.Connection.tangentSMul_mem_expDomain,
    LocallyConvex.Manifolds.Connection.starConvex_expDomainAt}
  \leanok
  \uses{thm:geodesic-spray-flow, def:exp, thm:maximal-flow}
  Let $M$ be Hausdorff and modeled on a Banach space, with a connection. The domain $D(\exp)$ is
  open in $\T M$, contains the zero section and is star-shaped in the fibers, and
  $\exp\colon D(\exp) \to M$ is smooth. For each $p \in M$, $D(\exp_p)$ is an open star-shaped
  neighbourhood of $0$ in $T_pM$, and $\exp_p$ is smooth on it.
\end{theorem}

\begin{proof}
  \leanok
  By Theorem~\ref{thm:geodesic-spray-flow}, $D(\exp) = \{v : (1, v) \in D(\mathrm{Fl}_S)\}$ and
  $\exp = \pi \circ \mathrm{Fl}_S(1, \cdot)$; the domain of the maximal flow is open and the flow
  is smooth on it (Theorem~\ref{thm:maximal-flow}, applied to the Hausdorff manifold $\T M$).
  Star-shapedness is $tv \in D(\exp) \iff t \in J_v$. The library works with smooth manifolds,
  so $\exp$ is smooth, rather than $C^{p-2}$ on a $C^p$ manifold as in \cite{Lan99}. That
  $\exp_p$ is a local diffeomorphism at $0$ is Theorem~\ref{thm:exp-local-diffeo}.
\end{proof}

\begin{theorem}[$\exp_p$ is a local diffeomorphism at $0$ {\cite[Thm.~IV.4.1, p.~108]{Lan99},
  \cite[Obs.~6.2.2, p.~155]{Lar22}}]
  \label{thm:exp-local-diffeo}
  \lean{LocallyConvex.Manifolds.exists_openPartialHomeomorph_of_contDiffOn,
    LocallyConvex.Manifolds.smoothMBMOn_trans_chart_symm,
    LocallyConvex.Manifolds.Connection.expAt_zero,
    LocallyConvex.Manifolds.Connection.contDiffOn_chart_expAt,
    LocallyConvex.Manifolds.Connection.hasFDerivAt_expAt_zero,
    LocallyConvex.Manifolds.Connection.exists_expAt_openPartialHomeomorph}
  \leanok
  \uses{thm:exp-smooth}
  Let $M$ be Hausdorff and modeled on a Banach space, with a connection, and let $p \in M$.
  Read in the preferred chart $\varphi$ at $p$, the derivative at $0$ of
  $w \mapsto \varphi(\exp_p w)$ is the identity: $d(\exp_p)(0_p) = \mathrm{id}$. Hence $\exp_p$
  restricts to a diffeomorphism from an open neighbourhood of $0$ in $T_pM$, contained in
  $D(\exp_p)$, onto an open neighbourhood of $p$.
\end{theorem}

\begin{proof}
  \leanok
  Since $\exp_p(tw) = \gamma_w(t)$ for $t \in J_w$ and $\gamma_w'(0) = w$, the derivative of
  $\varphi \circ \exp_p$ at $0$ in the direction $w$ is $w$. Michal--Bastiani smooth maps between
  open sets of Banach spaces are Fr\'echet smooth, so the inverse function theorem applies to
  $\varphi \circ \exp_p$; on the open set where its derivative is invertible the local inverse is
  smooth, and composing with $\varphi^{-1}$ gives the diffeomorphism onto an open set of $M$.
\end{proof}

\begin{theorem}[$G = (\pi, \exp)$ is a local diffeomorphism at $0_p$
  {\cite[Prop.~VIII.5.1, p.~215]{Lan99}, \cite[\S6.2.1, p.~156]{Lar22}}]
  \label{thm:exp-pair-local-diffeo}
  \lean{LocallyConvex.Manifolds.Connection.expChart,
    LocallyConvex.Manifolds.Connection.expChartDomain,
    LocallyConvex.Manifolds.Connection.contDiffOn_expChart,
    LocallyConvex.Manifolds.Connection.expChart_zero,
    LocallyConvex.Manifolds.prodShear,
    LocallyConvex.Manifolds.Connection.hasFDerivAt_expChart,
    LocallyConvex.Manifolds.Connection.exists_expChart_openPartialHomeomorph,
    LocallyConvex.Manifolds.Connection.exists_expPair_openPartialHomeomorph,
    LocallyConvex.Manifolds.Connection.exists_expPair_normal}
  \leanok
  \uses{thm:exp-local-diffeo}
  Under the hypotheses of Theorem~\ref{thm:exp-local-diffeo}, the map
  $G(v) = (\pi v, \exp v)$, $D(\exp) \to M \times M$, has differential
  $(a, b) \mapsto (a, a + b)$ at $0_p$ in charts, and restricts to a diffeomorphism from an open
  neighbourhood $V \subseteq D(\exp)$ of $0_p$ onto an open neighbourhood of $(p, p)$. There is an
  open $W \ni p$ with $W \times W \subseteq G(V)$: any $x, y \in W$ are joined by a unique geodesic
  $t \mapsto \exp(tv)$, $0 \le t \le 1$, with $v \in V$, and $v$ depends smoothly on $(x, y)$
  \cite[Teo.~6.2.3 (1)]{Lar22}, \cite[Cor.~VIII.5.2 (1)]{Lan99}.
\end{theorem}

\begin{proof}
  \leanok
  In the chart of $\T M$ over $\varphi$ and the chart $\varphi \times \varphi$ of $M \times M$,
  $G$ reads $(x, u) \mapsto (x, \varphi(\exp(x, u)))$. Since $G(x, 0) = (x, x)$ and
  $d(\exp_p)(0) = \mathrm{id}$ (Theorem~\ref{thm:exp-local-diffeo}), its differential at
  $(\varphi p, 0)$ is $(a, b) \mapsto (a, a + b)$, invertible; the inverse function theorem
  concludes. Uniqueness is injectivity of $G$ on $V$.
\end{proof}

\begin{theorem}[Uniformly normal neighbourhoods {\cite[Teo.~6.2.3 (2), p.~156]{Lar22},
  \cite[Cor.~VIII.5.2 (2), p.~216]{Lan99}}]
  \label{thm:uniformly-normal}
  \lean{LocallyConvex.Manifolds.OpenPartialHomeomorph.slice,
    LocallyConvex.Manifolds.OpenPartialHomeomorph.contDiffOn_slice,
    LocallyConvex.Manifolds.OpenPartialHomeomorph.contDiffOn_slice_symm,
    LocallyConvex.Manifolds.Connection.tangentChart_symm_proj_chart,
    LocallyConvex.Manifolds.Connection.exists_uniformlyNormal}
  \leanok
  \uses{thm:exp-pair-local-diffeo}
  Under the hypotheses of Theorem~\ref{thm:exp-local-diffeo}, let $\varphi$ be the preferred chart
  at $p$ and read tangent vectors in the chart $T$ of $\T M$ over $\varphi$. There are an open
  $W \ni p$ in the domain of $\varphi$ and $r > 0$ such that for every $q \in W$ the vectors of
  $T_qM$ with coordinates $(\varphi q, u)$, $\|u\| < r$, lie in $D(\exp)$, and
  $u \mapsto \exp_q(T^{-1}(\varphi q, u))$ is a diffeomorphism of $B_r(0) \subseteq E$ onto an open
  set containing $W$.
\end{theorem}

\begin{proof}
  \leanok
  The local inverse of $G$ in charts (Theorem~\ref{thm:exp-pair-local-diffeo}) is defined on a
  neighbourhood of $(\varphi p, 0)$ in $E \times E$, which contains a product of balls
  $B_r(\varphi p) \times B_r(0)$; its image is an open neighbourhood of $(\varphi p, \varphi p)$,
  which contains $B_\delta(\varphi p) \times B_\delta(\varphi p)$. Take
  $W = \varphi^{-1}(B_{\min(\delta, r)}(\varphi p))$. Since $G$ preserves the first coordinate,
  slicing at $x = \varphi q$ gives the diffeomorphism at $q$. No local compactness is used: the
  radius $r$ is uniform because it comes from the inverse function theorem applied to $G$ at the
  single point $0_p$. We do not claim this argument is new. The book also states the result for
  every $r \le R$ and that the joining geodesic stays in $\exp_x(B_R(0))$; these refinements are
  not formalized.
\end{proof}

\begin{theorem}[Complete Finsler manifolds are geodesically complete {\cite[Prop.~9.2.1]{Lar22}}]
  \label{thm:finsler-geodesically-complete}
  \lean{LocallyConvex.Manifolds.Connection.velocity_maxGeodesic_zero,
    LocallyConvex.Manifolds.Connection.not_bddAbove_maxInterval_of_complete,
    LocallyConvex.Manifolds.Connection.maxInterval_eq_univ_of_complete,
    LocallyConvex.Manifolds.Connection.expDomain_eq_univ_of_maxInterval_eq_univ,
    LocallyConvex.Manifolds.Connection.expDomainAt_eq_univ,
    LocallyConvex.Manifolds.Connection.exists_expDomainAt_eq_univ,
    LocallyConvex.Manifolds.Connection.expDomain_eq_univ_of_complete,
    LocallyConvex.Manifolds.Connection.exists_expDomainAt_eq_univ_of_complete}
  \leanok
  \uses{thm:finsler-geodesic-extension, thm:maximal-geodesic, def:exp}
  Under the hypotheses of Theorem~\ref{thm:finsler-geodesic-extension}, assume that every
  geodesic on a bounded open interval has bounded speed $\mu(\dot\alpha)$ and bounded backward
  speed $\mu(-\dot\alpha)$. Each assertion implies the next:
  \begin{enumerate}
    \item $M$ is complete for the symmetrized distance;
    \item every maximal geodesic is defined on $\mathbb{R}$;
    \item $M$ is geodesically complete: $D(\exp) = \T M$;
    \item there is $p \in M$ with $D(\exp_p) = T_pM$.
  \end{enumerate}
\end{theorem}

\begin{proof}
  \leanok
  $1 \Rightarrow 2$: if $b = \sup J_v$ were finite, Theorem~\ref{thm:finsler-geodesic-extension}
  would extend $\gamma_v$ past $b$, against maximality; boundedness below reduces to boundedness
  above through $J_{-v} = -J_v$. $2 \Rightarrow 3 \Rightarrow 4$ are immediate from
  Definition~\ref{def:exp}. We do not claim this argument is new.
\end{proof}

\begin{remark}[What fails in infinite dimensions]
  In finite dimensions the Hopf--Rinow theorem makes 1--4 equivalent for Riemannian (and
  reversible Finsler) manifolds, and adds that $\exp_p$ is onto. In infinite dimensions the last
  conclusion fails: Atkin constructs a complete Riemannian Hilbert manifold, the unit sphere of
  $\R \times \ell^2$ with the metric induced by $|\xi|^2 + \langle Tx, x\rangle$ for a suitable
  strictly positive operator $T$, whose two poles are joined by no geodesic
  \cite[(3), p.~262, and Lemma~6, p.~265]{Atk75}, so that $\exp_p$ is not onto (see also
  \cite[p.~288]{Eke78}). Even shortest curves
  need not exist: on the ellipsoid image of the unit sphere of a Hilbert space by
  $A = \mathrm{diag}(1, 1 + 2^{-1}, 1 + 2^{-2}, \dots)$, a complete Riemannian Hilbert manifold,
  $e_0$ and $-e_0$ are at distance $\pi$ but every curve joining them is strictly longer
  (McAlpin, as presented in \cite[Ejemplo~10.4.2, pp.~295--296]{Lar22}). No local compactness is
  used in the proof above: the extension rests on the uniform existence time of geodesics with
  bounded initial velocity in a chart and on the lower bound of $\mu$.
\end{remark}
